\documentclass{article}
\usepackage[T1]{fontenc}
\usepackage{lmodern}
\usepackage{graphicx}
\usepackage{amsmath,amssymb}
\usepackage[a4paper,margin=2cm]{geometry}
\usepackage[absolute,overlay]{textpos}
\setlength{\TPHorizModule}{1mm}
\setlength{\TPVertModule}{1mm}
\usepackage[indonesian]{babel}
\usepackage{hyperref}
\setlength{\emergencystretch}{2em}
% unit: Halaman 31

\begin{document}

% === Halaman 31 (python/native) ===

31

\# Bungkus kunci dalam format PKCS1\_OAEP untuk padding

    public\_key = PKCS1\_OAEP.new(key.publickey())

    private\_key = PKCS1\_OAEP.new(key)

    \# Tampilkan informasi kunci

    print("Modulus (n):", key.n)

    print("Kunci publik (e):", key.e)

    print("Kunci privat (d):", key.d)

    print("Bilangan prima p:", key.p)

    print("Bilangan prima q:", key.q)

    return public\_key, private\_key

def encrypt\_message(message, public\_key):    

    \#Enkripsi pesan dengan PKCS1\_OAEP    

    return public\_key.encrypt(message.encode())

def decrypt\_message(ciphertext, private\_key):   

    \#Dekripsi pesan dengan PKCS1\_OAEP    

    return private\_key.decrypt(ciphertext).decode()

\end{document}
